\section{Notation}\label{sec:app:notation}

A reference table, to be consulted rather than read. Each symbol is defined where it is introduced, and
this appendix collects the definitions in one place with the conventions that generated them.
Section~\ref{sec:app:sp:units} classifies the same symbols by units instead, and is the place to look
when a condition mixes the two.

The model carries two accounts at once: a value account in units of the final good, and a material
account in tonnes. Most of the notation exists to keep them apart, and the rules below are what make a
symbol readable without the sentence that introduced it.

\subsection{Rules}\label{sec:app:notation:rules}

\begin{enumerate}
\item \textbf{A superscript names a use or a process; a subscript names time.} $K^Y$ is capital in
production, $K^R$ capital in recycling, $C^N$ the cost of extraction, $W^Y$ the waste from production.
The period subscript is suppressed wherever timing carries no content and written wherever it does,
which in practice means the transition equations. The exogenous time argument of a primitive
($F$, $a$, $C^N$, $C^D$, $c^c$, $c^T$) is suppressed on the same terms.

\item \textbf{A stock is dated at the beginning of its period; a flow occurs during it.} A transition
equation therefore states what the period-$t$ flows leave behind for $t+1$. This is capital-style
dating and it fixes, once, which period a costate attaches to: the period in which the transition it
prices is chosen.

\item \textbf{A calligraphic capital is a function, a base or an aggregate}, never a quantity of the
same kind as its Roman counterpart. $\mathcal F$ is the technology whose value is $Y$; $\mathcal R$ is
the recycling technology whose value is $R^R$; $\mathcal W$ is the waste stock whose per-period outflow
is $H$; $\mathcal D$ and $\mathcal N$ are the two waste bases; $\mathcal G^{\varpi}$ and
$\mathcal G^{K}$ are the two recycling margins.

\item \textbf{$\lambda$ is the multiplier on the goods constraint, in utils per final good, and every
shadow price is stated net of it.} Dividing by $\lambda$ is what puts $p^S$, $p^X$, $p^P$,
$p^{\mathcal W}$, $p^M$, $p^W$, $\zeta$ and $\Psi$ into final goods per tonne; see
\eqref{eq:sp:costates-def}. Sign conventions are carried in the definitions, not in the symbols.

\item \textbf{$\Omega$ is the material intensity of the final good and $\omega^j$ the relative
intensity of use $j$}, so that $\Omega^j=\Omega\omega^j$ carries units and $\omega^j$ is a pure number.
The same split applies to the durables coefficient: $\Phi^I=\Omega\phi^I$. The nature-attributable
intensities $\Omega^{N,S}$ and $\Omega^{D,S}$ convert tonnage to tonnage and are pure numbers despite
the capital letter, because the mass they carry never passes through $Y$. The two families are not
comparable, and that is the distinction of Section~\ref{sec:sp:setup:materialaccounting}.

\item \textbf{A bar is a bound, and it is the same kind of object in every case.} $\bar R$ the material
floor, $\bar a$ the yield ceiling, $\bar X_{\max}$ the exhaustion point of discovery.

\item \textbf{A per-period rate lies in $(0,1)$}: $\delta$, $\mu$, $\theta$, $d^W$. A growth rate, an
elasticity and a discount rate do not, and never take those letters.

\item \textbf{One symbol, one meaning --- \emph{in this document}.} The paper states the same model in
continuous time and does not make all the same choices. Two symbols differ outright and must not be
carried across: $\zeta$ is the multiplier on the intensity condition here and the marginal value of the
final good there, and $q$ is the shadow price of installed capital in both but with different content,
$1-\Phi^I\left(1-\sigma\right)\left(p^W-p^M\right)$ here against $1+\Phi^Kp^M$ there. Where the two
documents correspond, the correspondence is stated rather than assumed.
\end{enumerate}

\subsection{The symbols}\label{sec:app:notation:table}

\begin{center}
\begin{tabular}{llll}
\hline
Symbol & Meaning & Units & Defined in \\
\hline
\multicolumn{4}{l}{\textit{States}}\\
$K$ & capital, split $K=K^Y+K^R$ & final goods & \eqref{eq:sp:setup:capital-allocation}\\
$S$ & known reserves & tonnes & \eqref{eq:sp:setup:reserves-motion}\\
$X$ & cumulative discovery, ceiling $\bar X_{\max}$ & tonnes & \eqref{eq:sp:setup:reserves-motion}\\
$P$ & stock of pollution & tonnes & \eqref{eq:sp:setup:pollution-motion}\\
$M^K$ & material embodied in the capital stock & tonnes & \eqref{eq:sp:states-motion}\\
$\mathcal W$ & anthropogenic waste stock & tonnes & \eqref{eq:sp:setup:wastestock}\\
\multicolumn{4}{l}{\textit{Controls and flows}}\\
$C$ & consumption & final goods & \eqref{eq:sp:setup:utility}\\
$I$, $G$ & net and gross capital formation, $G=I+\delta K$ & final goods & \eqref{eq:sp:setup:reserves-motion}\\
$Y$ & output of the final good & final goods & \eqref{eq:sp:setup:output}\\
$N$, $D$ & extraction and discovery & tonnes & \eqref{eq:sp:setup:reserves-motion}\\
$R$ & raw material input, $R=N+R^R$ & tonnes & \eqref{eq:sp:setup:materials}\\
$R^e$ & effective material input, $R-\bar R$ & tonnes & \eqref{eq:sp:setup:output}\\
$R^R$ & recycled material & tonnes & \eqref{eq:sp:setup:recycling}\\
$W$ & waste generated & tonnes & \eqref{eq:sp:ledger-streams}\\
$H$ & handled flow, $\mu\mathcal W$ & tonnes & \eqref{eq:sp:setup:wastestock}\\
$T$ & treated waste, $\varpi H$ & tonnes & \eqref{eq:sp:setup:recycling-crs}\\
$\varpi$ & treated share of the handled flow & dimensionless & \eqref{eq:sp:setup:waste-cost}\\
$\Xi$ & waste reaching the environment & tonnes & \eqref{eq:sp:setup:Xi}\\
$\Omega$ & material intensity of the final good & tonnes/final good & \eqref{eq:sp:intensity}\\
\multicolumn{4}{l}{\textit{Technology and costs}}\\
$\mathcal F$ & final-goods technology & --- & \eqref{eq:sp:setup:output}\\
$a(x)$ & recycling yield, $x=K^R/T$ & dimensionless & \eqref{eq:sp:setup:recycling}\\
$\bar a$ & yield ceiling; $\bar a<1$ hard (H), $\bar a=1$ soft (S) & dimensionless & \eqref{eq:sp:setup:recycling}\\
$\varepsilon$ & elasticity of the yield function, $a'x/a$ & dimensionless & \eqref{eq:sp:setup:recycling-crs}\\
$\alpha$ & marginal recycling yield, $a-a'x=a(1-\varepsilon)$ & dimensionless & \eqref{eq:sp:setup:recycling-derivatives}\\
$\bar R$ & material floor & tonnes & \eqref{eq:sp:setup:output}\\
$C^N$, $C^D$ & extraction and exploration cost & final goods & \eqref{eq:sp:setup:extraction-cost}\\
$C^W$, $c^W$ & waste cost, total and per handled tonne & final goods & \eqref{eq:sp:setup:waste-cost}\\
$c^c$, $c^T$ & collection charge, treatment cost & final goods/tonne & \eqref{eq:sp:setup:waste-cost}\\
\multicolumn{4}{l}{\textit{Materials accounting}}\\
$\Phi^x$ & material embodied per unit of $x$ & tonnes/final good & \S\ref{sec:sp:setup:materialaccounting}\\
$\Omega^x$, $\omega^x$ & waste per unit of $x$, absolute and relative & see rule 5 & \S\ref{sec:sp:setup:materialaccounting}\\
$\phi^I$, $\Phi^I$ & durables coefficient, $\Phi^I=\Omega\phi^I$ & see rule 5 & \eqref{eq:sp:setup:ybalance_OmegaPhi}\\
$M^x$, $W^x$ & embodied material and waste from $x$ & tonnes & \S\ref{sec:sp:setup:materialaccounting}\\
$\mathcal D$, $\mathcal N$ & waste bases: drawn through $R$, displaced from nature & see rule 5 & \eqref{eq:sp:setup:base}\\
$\sigma$ & storage share, $\Phi^IG/R$ & dimensionless & \eqref{eq:sp:zeta-explicit}\\
\multicolumn{4}{l}{\textit{Prices and multipliers}}\\
$\lambda$ & multiplier on the goods constraint & utils/final good & \eqref{eq:sp:costates-def}\\
$p^S$, $p^X$, $p^P$ & shadow prices of reserves, discovery, pollution & final goods/tonne & \eqref{eq:sp:costates-def}\\
$p^{\mathcal W}$ & shadow price of the waste stock & final goods/tonne & \eqref{eq:sp:costates-def}\\
$p^W$ & shadow price of waste generation, $p^W=-p^{\mathcal W}$ & final goods/tonne & \eqref{eq:sp:sum:W}\\
$p^M$ & shadow price of embodied material & final goods/tonne & \eqref{eq:sp:costates-def}\\
$q$ & shadow price of installed capital & dimensionless & \eqref{eq:sp:sum:I}\\
$\zeta$ & multiplier on the intensity condition & final goods/tonne & \eqref{eq:sp:zeta-explicit}\\
$\Psi$ & value of a tonne entering production & final goods/tonne & \eqref{eq:sp:Psi}\\
$h$ & handling dividend of a stockpiled tonne & final goods/tonne & \S\ref{sec:app:sp}\\
$\mathcal G^{\varpi}$, $\mathcal G^{K}$ & gains from treatment and from recycling capital & final goods/tonne & \eqref{eq:sp:sum:gains}\\
\multicolumn{4}{l}{\textit{Preferences and rates}}\\
$\rho$, $\beta$ & time preference, discount factor $1/(1+\rho)$ & rate & \eqref{eq:sp:setup:utility}\\
$\delta$ & depreciation & rate & \eqref{eq:sp:setup:resource-constraint}\\
$\mu$ & handling share of the waste stock & rate & \eqref{eq:sp:setup:wastestock}\\
$\theta(P)$ & decay of the pollution stock & rate & \eqref{eq:sp:setup:pollution-motion}\\
$d^W$ & leakage of the treated stream & rate & \eqref{eq:sp:setup:Xi}\\
\hline
\end{tabular}
\end{center}

\subsection{The source}\label{sec:app:notation:source}

The Julia implementation uses ASCII identifiers rather than these symbols, deliberately: a stray
re-encoding on Windows silently corrupts Unicode operators in Julia source, and the corruption survives
a passing test suite. The correspondence is therefore a lookup, and the table in
\texttt{model/SYMBOLS.md} is the one place it is written down. \texttt{code\_style.md} fixes the
transliteration rules and records the two clashes the source has to break that this document does not:
$\mu$ (the handling share, against the returns exponent of the quantitative technology) and $\beta$
(the discount factor, against the CES capital share).
